%% PLANTILLA PARA ARTICULOS
\documentclass[10pt,twoside,spanish,a4paper]{article}
%%%%%%%%%%%%%%%%%%%%%%%%%%%%% no tocar
% Alfredo Cuesta. 01/01/2013. acuestai@ucm.es

%si se quiere usar ARIAL hay que poner estas 2 lineas:
%\usepackage[scaled]{uarial} %DESCOMENTAR
\renewcommand*\familydefault{\sfdefault}
%

\usepackage[compact]{titlesec}
 \titlespacing\section{0pt}{12pt plus 0pt minus 0pt}{0pt plus 0pt minus 0pt}
 \titlespacing\subsection{0pt}{5pt plus 0pt minus 0pt}{0pt plus 0pt minus 0pt}
 \titlespacing\subsubsection{0pt}{0pt plus 0pt minus 0pt}{0pt plus 0pt minus 0pt}

\usepackage[utf8]{inputenc}
\usepackage{
amssymb,
graphicx,
color,
url
%wrapfig %Inclusión de gr·ficos al lado de texto
}
\definecolor{olivegreen}{rgb}{0.0,0.5,0.0}

\usepackage{listings}
\lstset{language=Matlab,%
    basicstyle=\footnotesize,
    keywordstyle=\color{blue},%
    commentstyle=\color{olivegreen}\textit,%
    emph=[1]{ans},emphstyle=[1]\color{black}, %some words to emphasise
    emph=[2]{resample,timeseries,syms},emphstyle=[2]\color{blue}, %some words to emphasise
}

\linespread{1.3}%separacion de lineas
\parindent=0.0in%indentancion de parrafo
\parskip=6pt%separacion entre el parrafo siguiente
\setlength{\topmargin}{-0.5in}
\setlength{\oddsidemargin}{0in}
\setlength{\evensidemargin}{0in}
\setlength{\textwidth}{16cm}%{4.72in}
\setlength{\textheight}{22cm}%{7.08in}

\usepackage{babel}
\addto\shorthandsspanish{\spanishdeactivate{~<>}}
% correct bad hyphenation here
\hyphenation{nu-es-tro co-no-ci-mien-to par-ti-cu-lar ca-rac-te-ris-ti-cas ni-vel se-gui-mien-to que-re-mos ca-ta-lo-go}
\date{}


% --- customization ---

\renewcommand\labelitemi{$\bullet$}
\let\olditemize\itemize
\renewcommand{\itemize}{
  \olditemize
  \vspace{-0pt}
  \setlength{\itemsep}{6pt}
  \setlength{\parskip}{0pt}
  \setlength{\parsep}{0pt}
  %\setlength{\topsep}{-15pt}
}

\let\oldenumerate\enumerate
\renewcommand{\enumerate}{
  \oldenumerate
  \vspace{-2pt}
  \setlength{\itemsep}{1pt}
  \setlength{\parskip}{0pt}
  \setlength{\parsep}{0pt}
  %\setlength{\topsep}{-15pt}
}

\usepackage{enumitem}
\usepackage{verbatim}
\usepackage{fancyhdr}
\usepackage{multirow}
\usepackage{colortbl}
\usepackage{verbatim}
\usepackage{amsmath}

\begin{document}

\renewcommand{\tablename}{Tabla}
\renewcommand{\figurename}{Figura}
\title{
  \vspace{-0.2cm}
	\begin{center}
		\large
		\begin{tabular}{rrl}
		\multirow{4}{*}{\includegraphics[width=1.7cm]{UCM.eps} \hspace{1cm}}&Grado: & Ingeniería Electrónica de Comunicaciones\\
		&Asignatura:& Procesamiento de Señales \\
		&Curso:     & 2025/2026 \\
		\end{tabular}
	\end{center}
	\rule{\textwidth}{0.4pt}
	%
	\vspace{-1.5cm}
	\begin{center}
		\color{blue} \Huge
		\begin{tabular}{c}
	        \textbf{Práctica 6}\\
		\textbf{Variables aleatorias}\\
		\end{tabular}
	\end{center}
	\vspace{-0.5cm}
	\rule{\textwidth}{0.4pt}
}
\maketitle
 \thispagestyle{empty}
\pagestyle{fancy}

\lfoot{\includegraphics[width=0.8cm]{UCM.eps}}
\cfoot{\vspace{-0.6cm} Ingeniería Electrónica de Comunicaciones }
\rfoot{\vspace{-0.6cm} \thepage}
\lhead{Práctica}
\rhead{Práctica 6: Variables aleatorias}
\renewcommand{\headrule}{ \vspace{-0.2cm} \rule{\textwidth}{0.4pt}}
\renewcommand{\footrule}{\vspace{0.5cm} \rule{\textwidth}{0.4pt}\\ \vspace{-0.3cm} \\}

%**************************************************************<- para indicar pregunta de la pr·ctica
\newcounter{ctPreg}
\setcounter{ctPreg}{1}
\newcommand{\itemP}{ \item[{\color{blue}\fbox{\textbf{Tarea \arabic{ctPreg}:}}}] \addtocounter{ctPreg}{1}}  %
\newcommand{\itemB}{ \item[{\color{blue}$\bullet$} ]    }
\newcommand{\itemR}{ \item[{\color{blue}$\triangleright$} ]    }       %
%**************************************************************

\vspace{-8em}

\section*{Objetivos}

El objetivo de esta práctica es:
\begin{itemize}
	\item Utilizar el método de la ruleta para generar muestras de una variable aleatoria. 
	\item Caracterizar comportamiento de diferentes variables aleatorias. 
\end{itemize}

\section*{Generación de muestras de una Variable Aleatoria}
La generación de muestras de una variable aleatoria es importante para muchas aplicaciones prácticas, por ejemplo para simular procesos aleatorios que describen un fenómeno a través de una distribución de probabilidad,
o para resolver integrales de las que no se conoce una solución cerrada. %En particular, se conoce con el nombre de Monte-Carlo
Generalmente se requiere generar muestras (números aleatorios) de forma que la frecuencia de ocurrencia se corresponda con 
una distribución de probabilidad que describe el comportamiento de la variable aleatoria. 

En la práctica, los computadores no son capaces de generar números realmente aleatorios, en su lugar utilizan lo que se denomina números pseudoaleatorios.
Estos son números que se generan de forma procedimental (es decir, a través de un algoritmo determinista) pero desde un punto de vista estadístico
es posible tratarlos como si fueran aleatorios. 

Un algoritmo sencillo y muy utilizado para generar números pseudoaleatorios es el generador lineal congruente, cuyo funcionamiento se basa 
en generar una secuencia de números a través de una relación de recurrencia de la forma $X(n+1)=(a X_n + c)~\textrm{mod}~m$. Según la elección particular
de los parámetros $a$, $c$ y $m$, la secuencia generada presentará un comportamiento más o menos adecuado como secuencia aleatoria. Es posible
realizar diferentes pruebas que verifican si el algoritmo genera números con las propiedades estadísticas deseadas y así poder tratarlos como 
aleatorios.

\newpage
\subsection*{Generación de Variables Aleatorias en Matlab}
Matlab proporciona diferentes funciones para generar muestras de variables aleatorias con las distribuciones de probabilidad más comunes: 
\begin{itemize}
	\item \texttt{randi(N,P,Q)} - Genera una matriz de tamaño $PxQ$ con muestras de una variable discreta con N casos equiprobables: $p(x_i)=\frac{1}{N},i=0,1,...,N-1$.
	\item \texttt{rand(P,Q)} - Genera una matriz de tamaño $PxQ$ con muestras de una variable discreta con una distribución uniforme
	en el intervalo $[0,1]$: $p(x)=\mathcal{U}_{[0,1]}(x)$. A partir de esta se puede generar cualquier distribución uniforme. Dada $X\sim p_X(x)=\mathcal{U}_{[0,1]}(x)$,
	podemos construir una distribución uniforme arbitraria mediante $Y=a+(b-a)X$. Se demuestra trivialmente que $p_Y(x)\sim U_{[a,b]}(x)$.
	\item \texttt{randn(P,Q)} - Genera una matriz de tamaño $PxQ$ con muestras de una variable discreta con una distribución gaussiana
	de media $\mu=0$ y varianza $\sigma^2=1$: $p(x)=\mathcal{N}_{(0,1)}(x)$. Al igual que en el caso anterior, si $X\sim p_X(x)=\mathcal{N}_{[0,1]}(x)$,
	y hacemos $Y=\mu+\sigma X$, es fácil ver que $p_Y(x)\sim \mathcal{N}_{[\mu,\sigma]}(x)$.
\end{itemize}

Para generar otras funciones de probabilidad, debemos recurrir a los algoritmos que hemos visto en las clases teóricas, que presentamos 
brevemente a continuación.

\subsection*{Método de la ruleta}
Este método se inspira en el conocido juego de azar\footnote{Por si hay algún despistado, el juego consiste en lanzar una bola dentro de una
ruleta que gira, dividida en diferentes sectores. La bola puede caer en cualquiera de las divisiones. El juego en sí no tiene mucha gracia, así 
que para hacerlo más divertido se suelen apostar grandes cantidades de dinero.}. Primero se divide la ruleta en diferentes sectores, proporcionales
a las probabilidades de la función que se desea muestrear. A continuación se simula el lanzamiento, y se genera la muestra según en que sector
caiga la bola virtual. Para implementar en Matlab, la ruleta se hace corresponder con el intervalo $[0,1]$, que se divide en tramos proporcionales
a las probabilidades que se quiere muestrear, y el lanzamiento de la bola se simula mediante el muestreo de una v.a. uniforme en el mismo
intervalo, $V\sim U_{[0,1]}(x)$. El código se lista a continuación:

\begin{lstlisting}
function x = ruleta(px, N)
  Px = cumsum(px);
  v = rand(N,1);
  x = sum((v > Px)') + 1;
end
\end{lstlisting}

% \section*{Método de Box-Muller}
% Este método se utiliza para generar variables aleatorias aleatorias gaussianas a partir de uniformes. Se muestrean dos v.a. uniformes, $y_1,y_2\sim \mathcal{U}_{[0,1]}(x)$, a continuación se realiza la transformación\\
% $x_1=\sqrt{-2\textrm{log}(y_1)}cos(2\pi y_2)$,\\ $x_2=\sqrt{-2\textrm{log}(y_1)}sin(2\pi y_2)$.

% $y_1=\mathrm{e}^{-\frac{x_1^2+x_2^2}{2}}$
% $y_2=\frac{1}{2\pi}\mathrm{atan}(\frac{x_2}{x_1})$\\

% $\mid\frac{\partial (y_1, y_2)}{\partial (x_1, x_2)}\mid = 
% \left|\begin{matrix}
% 	-x_1 \mathrm{e}^{-\frac{x_1^2+x_2^2}{2}} & -x_2 \mathrm{e}^{-\frac{x_1^2+x_2^2}{2}} \\
% 	-\frac{1}{2\pi}\frac{x_2/x_1^2}{1 + (\frac{x_2}{x_1})^2} & \frac{1}{2\pi}\frac{1/x_1}{1 + (\frac{x_2}{x_1})^2}  
% \end{matrix}\right|=-\frac{1}{\sqrt{2\pi}}\mathrm{e}^{-\frac{x_1^2}{2}}\frac{1}{\sqrt{2\pi}}\mathrm{e}^{-\frac{x_2^2}{2}}$
% function x = BoxMuller(px, N)
%   Px = cumsum(px);
%   v = rand(N,1);
%   x = sum((v > Px)') + 1;
% end

\newpage
\subsection*{Método de muestreo por rechazo}
Este método parte de una idea parecida al de la ruleta. En esencia, se genera una muestra de un distribución conocida (p.ej. la uniforme) y
se utiliza una ruleta para aceptar o rechazar la muestra, con una probabilidad determinada según la proporción entre la distribución objetivo y la auxiliar.

\begin{lstlisting}
function x_aceptadas = muestreo_rechazo(N)
  M = 2; 
  x_aceptadas = [];
  N_aceptadas = 0;
  while N_aceptadas < N
    x_q = rand(N - N_aceptadas, 1);
    y = rand(N - N_aceptadas, 1);
    l = find(y*M.*q(x_q) - p(x_q) < 0); % Muestras aceptadas
    x_aceptadas = [x_aceptadas; x_q(l)];
    N_aceptadas = length(x_aceptadas); 
  end
end

function y = p(x)
  % Densidad probabilidad objetivo
  y = ...
end

function y = q(x)
  % Densidad probabilidad auxiliar
  y = 1
end
\end{lstlisting}
	
\subsection*{Muestreo de una mezcla}
Este método se utiliza para generar muestras de una v.a. con una densidad de probabilidad de la forma $p(x) = \sum_{i=1}^{N}p(x|x_i)p(x_i)$, es decir,
una mezcla de varias conocidas. La idea es utilizar una ruleta para decidir a cual de los grupos pertenece la muestra y a continuación utilizar
el método que corresponda según la distribucion en particular. A continuación se lista el pseudocódigo del algoritmo:
\begin{lstlisting}
N = 1000;
px = [0.4, 0.2, 0.4];
for i = 1:N
  tipo = ruleta(px, 1);
  x(i) = generar_muestra(tipo);
end

function x = generar_muestra(tipo)
  switch tipo,
    case 1, x = ...; break;
    ...
    case n, x = ...; break;
  end
end
\end{lstlisting}

\subsection*{Histograma}
El histograma es una representación en forma de diagrama de barras de las frecuencias de ocurrencia de los valores de una v.a. 
Para variables discretas, el histograma es simplemente un recuento de las muestras. Si la variable es continua, se agrupan los valores 
por intervalos, y se realiza el recuento de las muestras que caen en cada intervalo. Si normalizamos el histograma para que su 
area se 1, lo podemos interpretar directamente como una aproximación de la densidad de probabilidad de la v.a.
La normalización se consigue dividiendo por $N$, el número de muestras, si la v.a. es discreta. Si la variable es continua,
hay que dividir por $N\Delta x$, donde $\Delta x$ es la anchura del intervalo utilizado para el recuento.

Para calcular el histograma en Matlab tenemos la función \texttt{histogram}. Si ignoramos los valores devueltos, esta función muestra 
una gráfica con el histograma generado. 
\begin{lstlisting}
h = histogram(x);
\end{lstlisting}
Por el contrario, si recogemos el valor devuelto en una variable, la gráfica no se mostrará y a cambio tendremos disponible un objeto con
toda la información:
\begin{lstlisting}
>> h = histogram(x);
  h = 

  Histogram with properties:

             Data: [1x500 double]
           Values: [1x14 double]
          NumBins: 14
         BinEdges: [1x15 double]
         BinWidth: 0.0100
        BinLimits: [0.6000 0.7400]
    Normalization: 'count'
        FaceColor: 'auto'
        EdgeColor: [0 0 0]

  Show all properties
\end{lstlisting}
Podemos utilizar esta información para normalizar el histograma.
\begin{lstlisting}
h_normalizado = h.Values / h.BinWidth / N;
\end{lstlisting}

\newpage
\section*{Estadísticos}
Es sencillo calcular los estadísticos muestrales en Matlab directamente a partir de las fórmulas:
\begin{lstlisting}
% x contiene las muestras de la v.a.
N = length(x);
media = sum(x) / N;
vcm = (x * x') / N;
varianza = sum((x - media).^2) / N;
desviacion = sqrt(varianza); % sqrt(sum((x - media).^2) / N);
\end{lstlisting}

También es posible usar las funciones \texttt{mean}, \texttt{std} y \texttt{var}.
\begin{lstlisting}
media = mean(x);
varianza = var(x, 1); % El 1 es para normalizar por N en lugar de (N-1)
desviacion = std(x, 1);
\end{lstlisting}

Conocida la distribución de probabilidad de la v.a., en algunos casos es posible caracterizarla por completo mediante de uno o varios 
estadísticos. Por ejemplo, una distribución normal se parametriza por la media y varianza, es decir, el conocimiento de estos
dos estadísticos permite reconstruir por completo la función de densidad de probabilidad. En otros casos el cálculo de 
los estadísticos no será suficiente.

\newpage

\section*{Generación y análisis de variables aleatorias}
	\begin{itemize}
		\itemP  \textbf{Muestreo de variables aleatorias.}
			El objetivo de esta tarea es generar diferentes conjuntos de muestras de una variable aleatoria discreta mediante el método de la ruleta.
			\itemR $p_X(x_i)=[0.1,0.25,0.15,0.35,0.1,0.05]$.
			
			Para comprobar que las muestras generadas siguen la función de probabilidad propuesta:		
			\begin{enumerate}[label=\alph*)]
				\item Genera tres conjuntos de muestras de diferente tamaño: $N=100$, $N=1000$, $N=10000$.
				\item Representa en una gráfica los valores de las muestras.
				\item Representa, en otra gráfica, los histogramas normalizados y la función de probabilidad propuesta.
				\item Compara los histogramas obtenidos con la función de probabilidad. Explica la influencia de N en el resultado.
				\item Calcula su media, valor cuadrático medio y varianza. Indica si la media y varianza caracterizan la función de probabilidad. 
			\end{enumerate}

			\hspace{0.2cm}\color{blue}Nota: \color{black} \'{U}nicamente se pueden usar los generadores aleatorios de Matlab  \color{blue}randi\color{black}, \color{blue}randn \color{black} y \color{blue}rand\color{black}.  Por lo tanto, se recomienda utilizar el método de muestreo más sencillo que, utilizando dichas funciones, genere las muestras propuestas.

		\itemP \textbf{Caracterización de conjuntos muestrales desconocidos}

			El objetivo de esta tarea es caracterizar conjuntos muestrales de origen desconocido a partir de su histograma.
			Los conjuntos muestrales que hay que analizar se generarán utilizando la función \color{blue}\texttt{DatosTarea2}\color{black}. Para cada conjunto muestral:
			\begin{enumerate}[label=\alph*)]
				\item Representa el histograma normalizado.
				\item Indica si la variable es continua o discreta, y cuál es su función de probabilidad original.
				\item Calcula los estadísticos habituales (media, valor cuadrático medio y varianza).
				\item ¿Los estadísticos sirven para parametrizar correctamente la función de probabilidad de las muestras?
			\end{enumerate}

	\end{itemize}

\end{document}
